\ifdefined\gkver\else\def\gkver{student}\fi
\documentclass[\gkver]{gaokaozhenti}
\begin{document}

\chapter{2002年全国春季理综}
%% paperType: 理综物理部分

\begin{enumerate}
\item
%% number: 16
%% paperName: 2002年全国春季理综第 16 题 6 分
%% typeId: 1
%% type: 单选题
%% chapter: 功和能
%% point: 机械能守恒定律
%% method: 无
%% score: 6
%% degree: 750
%% duplicateId: 0
%% body:
下列四个选项的图中，木块均在固定的斜面上运动，其中图A、B、C中的斜面是光滑的，图D中的斜面是粗糙的，图A、B中的$F$为木块所受的外力，方向如图中箭头所示，图A、B、D中的木块向下运动，图C中的木块向上运动，在这四个图所示的运动过程中机械能守恒的是\xzanswer{C}

\fourchoices[answer=C, ispicture=true, h=2.6cm]
{figs/16a.jpg}
{figs/16b.jpg}
{figs/16c.jpg}
{figs/16d.jpg}
%% answer: C
%% memo:
\memoanswer{
本题原卷及材料中未提供详解，待补充。
}

\item
%% number: 17
%% paperName: 2002年全国春季理综第 17 题 6 分
%% typeId: 1
%% type: 单选题
%% chapter: 电场
%% point: 静电感应
%% method: 无
%% score: 6
%% degree: 650
%% duplicateId: 0
%% body:
如图所示，在两个固定电荷$+q$和$-q$之间放入两个原来不带电的导体，1、2、3、4为导体上的四个点，在达到静电平衡后，各点的电势分别是$\varphi_{1}$、$\varphi_{2}$、$\varphi_{3}$、$\varphi_{4}$，则\xzanswer{B}

\onepicture[w=6.5cm]{figs/17.jpg}

\fourchoices[answer=B]
{$\varphi_{4}>\varphi_{3}>\varphi_{2}>\varphi_{1}$}
{$\varphi_{4}=\varphi_{3}>\varphi_{2}=\varphi_{1}$}
{$\varphi_{4}<\varphi_{3}<\varphi_{2}<\varphi_{1}$}
{$\varphi_{4}=\varphi_{3}<\varphi_{2}=\varphi_{1}$}
%% answer: B
%% memo:
\memoanswer{
本题原卷及材料中未提供详解，待补充。
}

\item
%% number: 18
%% paperName: 2002年全国春季理综第 18 题 6 分
%% typeId: 1
%% type: 单选题
%% chapter: 机械波
%% point: 波的图象
%% method: 无
%% score: 6
%% degree: 850
%% duplicateId: 0
%% body:
图中所示为一简谐横波在某一时刻的波形图，已知此时质点 A 正向上运动，如图中箭头所示，由此可断定此横波\xzanswer{C}

\onepicture[w=6.2cm]{figs/18.jpg}

\fourchoices[answer=C]
{向右传播，且此时质点 B 正向上运动}
{向右传播，且此时质点 C 正向下运动}
{向左传播，且此时质点 D 正向上运动}
{向左传播，且此时质点 E 正向下运动}
%% answer: C
%% memo:
\memoanswer{
本题原卷及材料中未提供详解，待补充。
}

\item
%% number: 19
%% paperName: 2002年全国春季理综第 19 题 6 分
%% typeId: 1
%% type: 单选题
%% chapter: 原子和原子核
%% point: 玻尔模型
%% method: 无
%% score: 6
%% degree: 650
%% duplicateId: 0
%% body:
氢原子的能级是氢原子处于各个定态时的能量值，它包括氢原子系统的电势能和电子在轨道上运动的动能，氢原子的电子由外层轨道跃迁到内层轨道时，\xzanswer{A}
\fourchoices[answer=A]
{氢原子的能量减小，电子的动能增加}
{氢原子的能量增加，电子的动能增加}
{氢原子的能量减小，电子的动能减小}
{氢原子的能量增加，电子的动能减小}
%% answer: A
%% memo:
\memoanswer{
本题原卷及材料中未提供详解，待补充。
}

\item
%% number: 20
%% paperName: 2002年全国春季理综第 20 题 6 分
%% typeId: 1
%% type: 单选题
%% chapter: 运动和力
%% point: 动量守恒定律
%% method: 无
%% score: 6
%% degree: 650
%% duplicateId: 0
%% body:
在高速公路上发生一起交通事故，一辆质量为$1500\Ukg$向南行驶的长途客车迎面撞上了一质量为$3000\Ukg$向北行驶的卡车，碰后两车接在一起，并向南滑行了一小段距离后停止，根据测速仪的测定，长途客车碰前以$20\Ums$的速率行驶，由此可判断卡车碰前的行驶速率\xzanswer{A}
\fourchoices[answer=A]
{小于$10\Ums$}
{大于$10\Ums$小于$20\Ums$}
{大于$20\Ums$小于$30\Ums$}
{大于$30\Ums$小于$40\Ums$}
%% answer: A
%% memo:
\memoanswer{
本题原卷及材料中未提供详解，待补充。
}

\item
%% number: 21
%% paperName: 2002年全国春季理综第 21 题 6 分
%% typeId: 1
%% type: 单选题
%% chapter: 磁场
%% point: 洛伦兹力
%% method: 无
%% score: 6
%% degree: 850
%% duplicateId: 0
%% body:
图为云室中某粒子穿过铅板 P 前后的轨迹，室中匀强磁场的方向与轨迹所在平面垂直（图中垂直于纸面向里），由此可知此粒子\xzanswer{A}

\onepicture[w=6.0cm]{figs/21.jpg}

\fourchoices[answer=A]
{一定带正电}
{一定带负电}
{不带电}
{可能带正电，也可能带负电}
%% answer: A
%% memo:
\memoanswer{
本题原卷及材料中未提供详解，待补充。
}

\item
%% number: 22
%% paperName: 2002年全国春季理综第 22 题 6 分
%% typeId: 1
%% type: 单选题
%% chapter: 光学
%% point: 透镜
%% method: 无
%% score: 6
%% degree: 650
%% duplicateId: 0
%% body:
图中L是凸透镜，$OO'$是它的主轴，AB是垂直于主轴的肖源，P是垂直于主轴的光屏，当两者到透镜的距离相等时，在光屏上得到清晰的像，如将AB向右移动任意一段距离后，再移动P，则在P上\xzanswer{D}

\onepicture[w=4.2cm]{figs/22.jpg}

\fourchoices[answer=D]
{总能得到缩小的像}
{总能得到较大的像}
{可能得到放大的像，也可能得到缩小的像}
{可能得到放大的像，也可能得不到像}
%% answer: D
%% memo:
\memoanswer{
本题原卷及材料中未提供详解，待补充。
}

\item
%% number: 23
%% paperName: 2002年全国春季理综第 23 题 6 分
%% typeId: 1
%% type: 单选题
%% chapter: 运动和力
%% point: 牛顿第二定律
%% method: 无
%% score: 6
%% degree: 550
%% duplicateId: 0
%% body:
质量为 $m$ 的三角形木楔 A 置于倾角为 $\theta$ 的固定斜面上，它与斜面间的动摩擦因数为 $\mu$，一水平力 $F$ 作用在木楔 A 的竖直平面上，在力 $F$ 的推动下，木楔 A 沿斜面以恒定的加速度 $a$ 向上滑动，则 $F$ 的大小为\xzanswer{C}

\onepicture[w=4.6cm]{figs/23.png}

\fourchoices[answer=C]
{$\frac{m[a + g(\sin \theta  + \mu \cos \theta )]}{\cos \theta }$}
{$\frac{m(a - g\sin \theta )}{\cos \theta  + \mu \sin \theta }$}
{$\frac{m[a + g(\sin \theta  + \mu \cos \theta )]}{\cos \theta  - \mu \sin \theta }$}
{$\frac{m[a + g(\sin \theta  + \mu \cos \theta )]}{\cos \theta  + \mu \sin \theta }$}
%% answer: C
%% memo:
\memoanswer{
本题原卷及材料中未提供详解，待补充。
}

\item
%% number: 24
%% paperName: 2002年全国春季理综第 24 题 6 分
%% typeId: 1
%% type: 单选题
%% chapter: 光学
%% point: 平面镜成像
%% method: 无
%% score: 6
%% degree: 450
%% duplicateId: 0
%% body:
一个点源 $S$ 对平面镜成像，设光源不动，平面镜以速率 $v$ 沿 $OS$ 方向向光源平移，镜面与 $OS$ 方向之间的夹角为 $30^\circ$，则光源的像 $S'$ 将\xzanswer{B}

\onepicture[w=7.0cm]{figs/24.jpg}

\fourchoices[answer=B]
{以速率 $0.5v$ 沿 $S'S$ 连线向 $S$ 运动}
{以速率 $v$ 沿 $S'S$ 连线向 $S$ 运动}
{以速率 $\sqrt{3}v$ 沿 $S'S$ 连线向 $S$ 运动}
{以速率 $2v$ 沿 $S'S$ 连线向 $S$ 运动}
%% answer: B
%% memo:
\memoanswer{
本题原卷及材料中未提供详解，待补充。
}

\item
%% number: 29
%% paperName: 2002年全国春季理综第 29 题 18 分
%% typeId: 4
%% type: 实验题
%% chapter: 电路
%% point: 伏安法测电阻
%% method: 无
%% score: 18
%% degree: 550
%% duplicateId: 0
%% body:
图甲、乙是两组同样的器材实物图，用来测量待测电阻 $R$ 的阻值，每组器材中包括：电池，电键，变阻器，电压表，电流表，等测电阻 $R$，若干导线。

\begin{enumerate}
	\item
	如果待测电阻 $R$ 的阻值比电压表的内阻不是小很多，但 $R$ 的阻值比电流表的内阻大很多，试在图甲中连线使之成为测量电路；如果待测电阻 $R$ 的阻值比电流表的内阻不是大很多，但 $R$ 的阻值比电压表的内阻小很多，试在图乙中连线使之成为测量电路。
	\item
	如果已知上述电压表的内阻 $R_{V}$ 和电流表的内阻 $R_{A}$，对图甲和图乙中连成的测量电路，分别写出计算待测电阻的公式（用测得的量和给出的电表内阻来表示）。
\end{enumerate}

\drawpicanswer{\onepicture[align=right, w=9.5cm]{figs/29.jpg}}{\onepicture[align=right, w=9.5cm]{figs/29_an.jpg}}
%% answer:
\jdanswer{
\begin{enumerate}
	\item
	如上图（答案）所示。
	\item
	用图甲线路时，$R=\frac{U-IR_{A}}{I}$；用图乙线路时，$R=\frac{U}{I-\frac{U}{R_{V}}}$。
\end{enumerate}
}
%% memo:
\memoanswer{
（1）连线见题图答案（图甲为电流表外接、图乙为电流表内接）。

（2）用图甲线路时，$R=\frac{U-IR_{A}}{I}$ ①；用图乙线路时，$R=\frac{U}{I-\frac{U}{R_{V}}}$ ②

评分标准：本题共 18 分，第（1）小题 10 分，第（2）小题 8 分。（1）图中有一处连接错误，不给分；变阻器按限流接法，只要连接正确，同样给分。（2）公式应与连线图对应，如果考生连线答案内、外接相反，但只要公式对应并正确，不重复扣分。
}

\item
%% number: 30
%% paperName: 2002年全国春季理综第 30 题 28 分
%% typeId: 6
%% type: 计算题
%% chapter: 气体性质
%% point: 玻意耳定律
%% method: 无
%% score: 28
%% degree: 550
%% duplicateId: 0
%% body:
如图所示，竖直放置的气缸内盛有气体，上面被一活塞盖住，活塞通过劲度系数 $k=600\UNm$ 的弹簧与气缸相连接，系统处于平衡状态，已知此时外界大气压强 $p_{0}=1.00\times10^{5}\UNmq$，活塞到缸底的距离 $l=0.500\Um$，缸内横截面积 $S=1.00\times10^{-2}\Umq$，今在等温条件下将活塞缓慢上提到距缸底为 $2l$ 处，此时提力为 $F=500\UN$，弹簧的原长 $l_{0}$ 应为多少？若提力为 $F=700\UN$，弹簧的原长 $l_{0}$ 又应为多少？

不计摩擦及活塞和弹簧的质量，并假定在整个过程中，气缸不漏气，弹簧都遵从胡克定律。
\onepicture[align=right, h=4.8cm]{figs/30.jpg}
%% answer:
\jdanswer{
\begin{enumerate}
	\item
	当 $F=500\UN$ 时，$l_{0}=1.50\Um$。
	\item
	当 $F=700\UN$ 时，$l_{0}=0.833\Um$。
\end{enumerate}
}
%% memo:
\memoanswer{
参考解答一：

设弹簧的原长为 $l_{0}$，气体原来压强为 $p$，后来为 $p'$，则由玻意耳定律可得

$pl=p'\cdot2l$ ①

在原来状态下，活塞受力如图 1 所示，由力学平衡可得

$pS=p_{0}S+k(l-l_{0})$ ②

在后来状态下，活塞受力如图 2 所示，由力学平衡可得

$p'S+F=p_{0}S+k(2l-l_{0})$ ③

由①、②、③联立解得

$p=\frac{2(F-kl)}{S}$ ④

由式得

$l_{0}=l+\frac{S}{k}(p_{0}-p)$ ⑤

当 $F=500\UN$ 时，由④式得 $p=0.4p_{0}$，再代入⑤式得 $l_{0}=1.50\Um$，可见在整个过程中弹簧始终处于压缩状态。

当 $F=700\UN$ 时，由④式得 $p=0.8p_{0}$，再代入⑤式得 $l_{0}=0.833\Um$，可见在过程开始时弹簧处于压缩状态，当活塞提高到距缸底距离超过 $l_{0}=0.833\Um$ 后，弹簧被拉伸。

评分标准：本题共 28 分。正确求出 $l_{0}$ 的两个结果。

参考解答二：

设开始时弹簧的压缩量为 $x$（当得出 $x$ 为负值时，表示开始时弹簧被拉长），原长为 $l_{0}$。

依题意得方程：

$p_{0}S=pS+kx$ ①

$p_{0}S=p'S-k(l_{0}-2x)+F$ ②

$p'S\cdot2(l_{0}-x)=pS(l_{0}-x)$ ③

$l_{0}=l+x$ ④

由①、②、③、④式联立，解得

$x=\frac{p_{0}S-2F+2kl}{k}$ ⑤

当 $F=500\UN$ 时，代入⑤式，得 $x=1.00\Um$，$l_{0}=1.50\Um$。

当 $F=700\UN$ 时，代入⑤式，得 $x=0.333\Um$，$l_{0}=0.833\Um$。

（用其他解法，正确列出公式，并正确求得 $l_{0}$ 值，同样给分。）

\twopicture[showlabel=false, h=2.0cm]{figs/30_an1.jpg}{figs/30_an2.jpg}
}

\item
%% number: 31
%% paperName: 2002年全国春季理综第 31 题 20 分
%% typeId: 6
%% type: 计算题
%% chapter: 电磁感应
%% point: 切割磁感线电动势
%% method: 无
%% score: 20
%% degree: 550
%% duplicateId: 0
%% body:
磁铁在电器中有广泛的应用，如发电机，如图所示，已知一台单和发电机转子导线框共有 $N$ 匝，线框长为 $l_{1}$，宽为 $l_{2}$，转子的转动角速度为 $\omega$，磁极间的磁感强度为 $B$，试导出发电机的瞬时电动势 $E$ 的表达式。

现在知道有一种强永磁材料钕铁硼，用它制成发电机的磁极时，磁感强度可增大到原来的 $k$ 倍，如果保持发电机结构和尺寸，转子转动角速度，需产生的电动热都不变，那么这时转子上的导线框需要多少匝？
\onepicture[align=right, h=4.8cm]{figs/31.jpg}
%% answer:
\jdanswer{
\begin{enumerate}
	\item
	$E_{N}=Nl_{1}l_{2}\omega B\sin\omega t$。
	\item
	$N'=\frac{N}{k}$。
\end{enumerate}
}
%% memo:
\memoanswer{
取轴 $Ox$ 垂直于磁感强度，线框转角为 $\theta$（如图所示），线框长边垂直于纸面，点 $A$、$B$ 表示线框长边导线与纸面的交点，$O$ 点表示转轴与纸面的交点。

线框长边的线速度为

$v=\frac{\omega l_{2}}{2}$ ①

一根长边导线产生的电动势为 $\frac{\omega l_{2}}{2}\cdot B\sin\theta\cdot l_{1}$，一匝导线框所产生的感应电动势为

$E_{1}=l_{1}l_{2}\omega B\sin\theta$ ②

$N$ 匝线框产生的电动势应为

$E_{N}=NE_{1}=Nl_{1}l_{2}\omega B\sin\omega t$ ③

磁极换成钕铁硼永磁体时，设匝数为 $N'$，则有

$E_{N'}=N'l_{1}l_{2}\omega kB\sin\omega t$ ④

由 $E_{N}=E_{N'}$ 可得

$N'=\frac{N}{k}$ ⑤

如果考生取磁场线方向为轴并得出正确结果同样给分。

\onepicture[align=right, h=3.6cm]{figs/31_an.jpg}
}

\end{enumerate}

\end{document}
